\section{Discussion}
\label{sec:discussion}

The design fingerprint offers something that TIGER profiles alone do not: a representation
of a game's \emph{relative} emphasis that is directly comparable across quality tiers.
A high-quality gateway game and a high-quality heavy euro can have very similar TIGER
coordinates but very different fingerprints if the gateway game's elegance is achieved
through reduction while the euro's elegance is achieved through integration.

\subsection{Implications for Designers}

For a designer, the fingerprint is a diagnostic: it identifies which axes the current
design is deprioritizing relative to its own average, flagging potential blind spots.
A game with a low-distinctiveness fingerprint ($\sigma_g < 0.5$) is not necessarily
bad, but it is generic: it scores near its own average on all axes, suggesting that
no design axis has been pushed hard enough to create a distinctive play experience.

\subsection{Fingerprints and Designer Identity}

The finding that games by prolific designers cluster in fingerprint space (if confirmed
empirically) has implications for understanding designer identity as a measurable
design-space trajectory rather than a stylistic label. This connection between fingerprint
analysis and computational creativity research on artistic style is a direction for
future work.

\subsection{Limitations}

Within-game z-scoring requires that each game has meaningful variation across axes.
Games where all axes are scored identically have undefined z-scores. More practically,
the fingerprint captures relative emphasis within the 32-axis framework, which is itself
a constructed sampling of design space --- different axis sets would produce different
fingerprints for the same game.
